\section*{Chapitre \ref{ch:qm:resampling} --- Rééchantillonnage}

\begin{solution}{exo:qm:resampling:1}
Chacun des $n$ tirages ne retient pas le jour avec la probabilité $1 - 1/n$, d'où $(1 - 1/n)^n$~: 0.3677 pour $n = 1\,260$, et $e^{-1} = 0.3679$ dans la
limite. Un rééchantillon contient environ 63~\% des jours distincts.
\end{solution}

\begin{solution}{exo:qm:resampling:2}
Le bloc se prolonge chaque jour avec la probabilité $0.95$, si bien qu'il dure plus de 60 jours avec la probabilité $0.95^{60} = 0.046$.
\end{solution}

\begin{solution}{exo:qm:resampling:3}
$5^5 = 3\,125$ rééchantillons ordonnés~; $\binom{9}{5} = 126$ multi-ensembles (choix de cinq tirages parmi cinq valeurs avec répétition).
\end{solution}

\begin{solution}{exo:qm:resampling:4}
Soit $S = \sum_j(X_j - \bar X)^2$ et $d_i = X_i - \bar X$. Retirer $X_i$ laisse la somme des carrés $S - \frac{n}{n-1}d_i^2$ par rapport à la nouvelle
moyenne, d'où $\hat\theta_{(i)} = \bigl(S - \frac{n}{n-1}d_i^2\bigr)/(n - 1)$. Comme $\sum_i d_i^2 = S$, la moyenne vaut $\bar\theta = (S - \frac{S}{n-1})/(n
- 1) = \frac{(n-2)S}{(n-1)^2}$. L'\omterm{def:qm:estimation:estimator}{estimateur} corrigé vaut $n\hat\theta - (n - 1)\bar\theta = S - \frac{(n-2)S}{n-1} = \frac{S}{n-1} = s^2$.
\end{solution}

\begin{solution}{exo:qm:resampling:5}
Avec $\gamma(k) = \gamma(0)\phi^{|k|}$~: $G = 2\gamma(0)\phi/(1 - \phi)^2$ et $g(0) = \gamma(0)(1 + \phi)/(1 - \phi)$, d'où $G/g(0) = 2\phi/(1 -
\phi^2) = 4/3$, et $b = (4/3)^{2/3} \times 1\,260^{1/3} = 1.211 \times 10.80 = 13.1$ jours.
\end{solution}

\begin{solution}{exo:qm:resampling:6}
$1 - 0.99^{100} = 0.634$. Le maximum \omterm{def:qm:resampling:bootstrap}{bootstrap} égale le maximum de l'échantillon dans 63~\% des cas~: sa loi présente un atome, alors que la loi vraie du
maximum est continue, si bien que le \omterm{def:qm:resampling:bootstrap}{bootstrap} est incohérent pour les extrêmes.
\end{solution}

\begin{solution}{exo:qm:resampling:7}
\texttt{spurious\_rejections()} avec 199 permutations ou décalages par paire~: le \omterm{def:qm:resampling:perm}{test de permutation} rejette dans 21.5~\% des 200 paires, soit quatre
fois sa taille nominale, car les deux séries sont persistantes et le mélange détruit cette propriété~; le test par décalages circulaires, qui conserve la dépendance propre à chaque
série, rejette dans 2.5~\%.
\end{solution}

\begin{solution}{exo:qm:resampling:8}
Trois défauts. Rééchantillonner les jours de façon indépendante ignore le regroupement du P\&L~; pour le vendeur de volatilité du chapitre, de tels intervalles contiennent
la vraie valeur dans 68.5~\% des cas, et non 95~\%. Le \omterm{def:qm:resampling:bootstrap}{bootstrap} ne peut pas produire de pertes supérieures à la pire journée de l'échantillon, si bien que même les intervalles par
blocs sont trop étroits pour une stratégie à pertes rares mais importantes. Le percentile est une affirmation sur les rééchantillons de cet historique, et non une
probabilité portant sur le vrai \omterm{def:qm:estimation:sharpe}{ratio de Sharpe}~; si la stratégie a été choisie parmi de nombreuses variantes, la correction du chapitre~12 s'ajoute.
\end{solution}

\begin{solution}{pb:qm:resampling:1}
\textbf{1.} \omterm{def:qm:estimation:sharpe}{Ratio de Sharpe} 1.11, asymétrie $-4.7$, aplatissement 34.
\textbf{2.} 0.13, 0.10 et $-0.02$.
\textbf{3.} La stratégie verse $r_t^2$, et les rendements au carré sont autocorrélés parce que la volatilité est persistante~: les journées calmes suivent les journées calmes,
les pertes suivent les pertes.
\textbf{4.} 3.9 fois la variance ($1 + 2\sum_k(e^{0.64 \cdot 0.9^k} - 1)/(3e^{0.64} - 1)$).
\textbf{5.} 0.52~: le terme d'asymétrie $-\gamma_3\mathrm{SR}$ s'ajoute à la variance, mais la formule suppose encore des jours indépendants.
\textbf{6.} $(0.17, 2.25)$, \omterm{def:qm:estimation:estimator}{erreur type} 0.53.
\textbf{7.} Longueur moyenne des blocs de 18.6 jours pour le \omterm{def:qm:resampling:blocks}{bootstrap stationnaire}, 21.2 jours pour le \omterm{def:qm:resampling:bootstrap}{bootstrap} par blocs circulaire.
\textbf{8.} $(-0.34, 2.93)$, \omterm{def:qm:estimation:estimator}{erreur type} 0.84.
\textbf{9.} $(-0.25, 2.96)$, \omterm{def:qm:estimation:estimator}{erreur type} 0.83.
\textbf{10.} \omterm{def:qm:estimation:estimator}{Erreur type} 0.53 (la réponse iid~: les suppressions d'un seul jour ignorent la dépendance) et biais 0.04, d'où un ratio corrigé de 1.06.
\textbf{11.} 1.10, à partir de 4\,000 historiques simulés.
\textbf{12.} Dans 68.5~\% des cas.
\textbf{13.} 82.75~\% avec la longueur de bloc automatique (\omterm{def:qm:resampling:blocks}{bootstrap par blocs mobiles} identique), 84.0~\% avec des blocs de 40 jours.
\textbf{14.} L'\omterm{def:qm:estimation:estimator}{erreur type} du \omterm{def:qm:resampling:bootstrap}{bootstrap} se stabilise autour de 0.97 face à une vraie valeur de 1.10~: les historiques sans épisode violent ne peuvent pas en rééchantillonner
un, et cinq ans contiennent peu de regroupements de volatilité indépendants.
\textbf{15.} Il contient la vraie valeur dans 91.5~\% des cas~; le reste du déficit tient aux queues épaisses du P\&L.
\textbf{16.} L'intervalle stationnaire, $(-0.25, 2.96)$, présenté comme une borne inférieure de l'incertitude, avec sa \omterm{def:qm:resampling:interval}{probabilité de couverture} mesurée sur un modèle de
la stratégie.
\textbf{17.} Non~: l'intervalle contient zéro, et l'incertitude réelle est encore plus grande~; un doublement reposerait sur l'intervalle iid, qui est le
mauvais.
\textbf{18.} Tests de résistance avec des régimes de volatilité absents de l'échantillon, l'historique de la stratégie à travers les pics de volatilité passés, un
backtest plus long, et la recherche qui a produit la stratégie.
\textbf{19.} Résultat nommé~: \emph{l'intervalle trop étroit}~: les intervalles \omterm{def:qm:resampling:bootstrap}{bootstrap} nominaux à 95~\% pour le \omterm{def:qm:estimation:sharpe}{ratio de Sharpe} du vendeur
de volatilité contiennent la vraie valeur dans 68.5~\% des cas lorsque les jours sont rééchantillonnés de façon indépendante, et dans 82.75~\% des cas
avec les blocs automatiques du \omterm{def:qm:resampling:blocks}{bootstrap stationnaire}.
\textbf{20.} Il doit préserver toute dépendance que l'\omterm{def:qm:testing-and-multiple-testing:test}{hypothèse nulle} ou l'estimande ne nie pas~: au sein d'une même journée, entre actifs, et
d'un jour à l'autre dans le temps.
\end{solution}

\begin{solution}{iq:qm:resampling:1}
Il estime la loi d'échantillonnage d'une statistique en la recalculant sur des rééchantillons tirés avec remise à partir des données, la loi empirique tenant
lieu de la vraie loi. Il échoue pour les statistiques gouvernées par les valeurs extrêmes (le maximum, les quantiles de queue), pour des données dépendantes rééchantillonnées comme si elles étaient
indépendantes, et pour de très petits échantillons.
\iqlookfor{Le principe de substitution, et au moins deux modes de défaillance.}
\end{solution}

\begin{solution}{iq:qm:resampling:2}
Le P\&L quotidien est dépendant~: la volatilité se regroupe, les positions se recouvrent, les edges se présentent par régimes. Le rééchantillonnage indépendant estime la variance de court terme,
et l'intervalle est trop étroit~; dans l'exemple du chapitre, il contient la vraie valeur dans 68.5~\% des cas au lieu de 95~\%. Rééchantillonner par blocs (\omterm{def:qm:resampling:blocks}{bootstrap stationnaire}), et vérifier
la \omterm{def:qm:resampling:interval}{probabilité de couverture} par simulation.
\iqlookfor{Variance de long terme et variance de court terme, rééchantillonnage par blocs, la couverture comme test.}
\end{solution}

\begin{solution}{iq:qm:resampling:3}
L'erreur de l'\omterm{def:qm:estimation:estimator}{estimateur} de la variance est minimisée par une longueur croissant comme $n^{1/3}$, la constante étant fixée par les autocovariances~; il faut l'estimer
(Politis--White avec la correction de 2009), puis vérifier la sensibilité en traçant le résultat en fonction de la longueur de bloc. Les données persistantes,
à queue épaisse, demandent des blocs plus longs que ne le suggère la règle.
\iqlookfor{La vitesse $n^{1/3}$, une règle automatique et un contrôle de sensibilité.}
\end{solution}

\begin{solution}{iq:qm:resampling:4}
Lorsque les séries sont sériellement dépendantes~: le mélange rend le signal échangeable alors qu'il ne l'est pas, et la corrélation fallacieuse entre deux
séries persistantes est bien plus fréquente que ne le montrent les mélanges. Dans le chapitre, un \omterm{def:qm:resampling:perm}{test de permutation} rejette 21.5~\% des paires indépendantes au
seuil de 5~\%. Utiliser des décalages circulaires ou des permutations par blocs.
\iqlookfor{L'\omterm{def:qm:resampling:perm}{échangeabilité} comme condition, et un schéma de rééchantillonnage qui conserve la dépendance.}
\end{solution}

\begin{solution}{iq:qm:resampling:5}
Le \omterm{def:qm:resampling:jackknife}{jackknife}~: $n$ recalculs déterministes, chacun privé d'une observation, une correction de biais qui élimine le terme en $1/n$, et une \omterm{def:qm:estimation:estimator}{erreur type} pour les statistiques
régulières. Le \omterm{def:qm:resampling:bootstrap}{bootstrap}~: une distribution complète, donc des intervalles et des tests, au prix de l'aléa et d'un calcul plus lourd. Les deux exigent
des blocs pour les données dépendantes.
\iqlookfor{Biais contre distribution, et besoin commun de blocs.}
\end{solution}

\begin{solution}{iq:qm:resampling:6}
Rééchantillonner conjointement les jours pour les deux cents stratégies (en conservant leur corrélation) avec un \omterm{def:qm:resampling:blocks}{bootstrap stationnaire} (en conservant la dépendance temporelle),
recentrer la moyenne rééchantillonnée de chaque stratégie sur sa moyenne empirique, calculer la statistique $t$ maximale par rééchantillon, et lire la $p$-valeur du
maximum observé à partir de cette distribution~: le \omterm{def:qm:testing-and-multiple-testing:rc}{Reality Check} de White, ou le SPA de Hansen avec studentisation et recentrage, ou la méthode de Romano--Wolf pour
variante descendante (step-down).
\iqlookfor{Rééchantillonnage conjoint, recentrage pour imposer l'\omterm{def:qm:testing-and-multiple-testing:test}{hypothèse nulle}, le maximum comme statistique.}
\end{solution}
